\subsection{Oracle-scale tail certificates without a known gap law}
\label{sec:arbitrary-law-tail}

The next result addresses unobserved small-gap mass. It is a paper proof,
with a tested executable controller, and is not part of the Lean audit.
The additional core allocation and null-likelihood argument needed for
oracle-rate estimation of the positive fraction is supplied in
\cref{sec:arbitrary-law-adaptation}.

Let $\mu$ be any probability law on $(0,1/2]$. A coin's sign may depend
arbitrarily on its persistent magnitude. Fresh coins are independent and
flips are conditionally independent given a coin's parameters. Only flips
are observed. Write $F(t)=\mu((0,t])$ and $\log_2^+x=\max(0,\log_2x)$.

\begin{theorem}[Arbitrary-law tail certificate]\label{thm:arbitrary-law-tail}
Fix $J\ge1$, $N=4^J$, $t=2^{-J}$ and $0<\delta<1$. Let $m$ be odd with
$m\ge8\log(2/\delta)$. There is an observable procedure using at most $mN$
physical flips that returns $U\in[0,1]$. Put
\[
 w_J(d)=\min\{N,1024d^{-2}(1+\log_2^+(d/t))\},\qquad
 W_J=\int w_J\,d\mu.
 \tag{AT1}
\]
If $u_J$ is the unresolved probability of the single-coin test defined below,
then
\[
 \tfrac12F(t)\le u_J\le\int\min\{1,16(t/d)^2\}\,\mu(dd).
 \tag{AT2}
\]
With probability at least $1-\delta$,
\[
 F(t)\le U\le\min\{1,128\max(u_J,W_J/N)\}.
 \tag{AT3}
\]
For each fixed $\mu$ and $0<\rho\le1/4$, there is a finite onset
$J_0(\mu,\rho)$ such that, at all $J\ge J_0$, the same event also gives
\[
 F(2^{-J})\le U\le524288\,\mathcal K_{\mu,\rho}(4^J).
 \tag{AT4}
\]
The coefficient is universal. The algorithm receives neither $\mu$, $\rho$
nor the onset. No inverse-gap moment assumption is made.
\end{theorem}

\paragraph{The single-coin test.}
Inspect a coin's mean after $n_k=4^k$ flips, $k=0,\ldots,J$, using
\[
 h_k=\left\lceil\sqrt{2n_k(J-k+2)}\right\rceil,
 \qquad r_k=h_k/n_k.
 \tag{AT5}
\]
Call it resolved and stop if $|\bar X_{n_k}-1/2|>r_k$; otherwise continue,
calling it unresolved if the terminal check fails. Let $T$ be its work and
$Y$ its unresolved indicator. Always $1\le T\le N$. Complementing every
flip preserves its entire stopping path, so the conditional laws of $T,Y$
depend only on the magnitude, including when sign and magnitude are dependent.

If $d\le t$, then $r_k\ge2/\sqrt{n_k}\ge2t\ge2d$. Resolving at check $k$
requires an empirical error larger than $r_k/2$. Two-sided Hoeffding gives
\[
 \Pr_d(Y=0)\le\sum_{k=0}^J2e^{-n_kr_k^2/2}
 \le\sum_{s=0}^{\infty}2e^{-(s+2)}
 =\frac{2}{e(e-1)}<\frac12.
 \tag{AT6}
\]
No independence between checks is used. At the terminal check $h_J=2\sqrt N$
exactly, so $r_J=2t$. If $d\ge4t$, being unresolved forces empirical error
at least $d/2$. Chebyshev and $\operatorname{Var}(\bar X_N)\le1/(4N)$ bound
its probability by $1/(Nd^2)$. For $d<4t$ use one. Integrating proves (AT2).

We also have $\mathbb E_dT\le w_J(d)$. Indeed,
$r_k\le3\sqrt{(J-k+2)/n_k}$. If it exists, let $k_0$ be the first check with
$n_{k_0}d^2\ge36(J-k_0+2)$. Then $k_0\ge1$, and for every $k\ge k_0$,
$r_k\le d/2$, giving $\Pr_d(T>n_k)\le2e^{-n_kd^2/2}$. Since
$e^{-x}\le2/x^2$ for $x>0$,
\[
 \mathbb E_dT
 \le n_{k_0}+6\sum_{k=k_0}^{J-1}n_ke^{-n_kd^2/2}
 \le n_{k_0}+\frac{64}{n_{k_0}d^4}\le2n_{k_0}.
 \tag{AT7}
\]
Here $n_{k_0}d^2\ge72$. Minimality gives
$n_{k_0}d^2<144(J-k_0+3)$, while $2^{k_0}d\ge1$ implies
$J-k_0\le\log_2(d/t)$. Thus the expected work is at most
$864d^{-2}(1+\log_2^+(d/t))$. If there is no such check, $Nd^2<72$ and
$T\le N<72d^{-2}$. Combining both cases with $T\le N$ proves (AT1)'s
work bound and $\mathbb ET\le W_J$.

\paragraph{Inverse sampling and the physical cap.}
Test fresh coins until the first unresolved coin or budget exhaustion.
Pause within a coin when necessary and count only completed tests. If $n$
coins have completed, return $V=\min(1,8/n)$, taking $V=1$ if $n=0$.
The first unresolved coin is included in $n$.

For fixed $J$, with probability at least $7/8$, this probe covers $F(t)$
simultaneously at every cumulative budget. Put $u=u_J$. When $F(t)>V$,
$n>8/F(t)\ge4/u$. Reaching such an $n$ requires the first
$\lfloor4/u\rfloor$ coins to be resolved, with probability at most
\[
 (1-u)^{\lfloor4/u\rfloor}\le e^{-4+u}\le e^{-3}<1/8.
 \tag{AT8}
\]
If $u=0$, (AT2) makes failure impossible. Thus random completed sample sizes
and observation-dependent budgets are covered by one event.

For a deterministic budget $b>0$, separately,
\[
 \Pr\{V\le\min(1,128\max(u,W_J/b))\}\ge3/4.
 \tag{AT9}
\]
Set $A=\max(u,W_J/b)>0$. If $A\ge1/16$ this is automatic. Otherwise let
$q=\lfloor1/(8A)\rfloor\ge1/(16A)$. In a hypothetical stream of $q$ full
tests, the chance any is unresolved is at most $qu\le1/8$. Markov bounds the
chance their combined work exceeds $b$ by $qW_J/b\le1/8$. Outside this union,
the actual probe completes at least $q$ tests, hence $V\le8/q\le128A$.
The argument charges partial tests too and does not assume independence of
the work and unresolved events.

Take the median of $m$ independent probes, each capped at $b=N$. A block's
validity failure has probability at most $1/8$ and its performance failure
at most $1/4$. Hoeffding across blocks bounds each median failure by
$e^{-m/8}$; their union is at most $\delta$. This proves (AT3), with the
physical cap valid on every path.

\paragraph{The oracle comparison.}
Let $\nu$ be the lowest-gap quantile submeasure of mass $\rho$, splitting a
boundary atom when needed, and set $\sigma=\mu-\nu$. There is $a>0$ such
that $\sigma$ is supported on $[a,1/2]$. The existing quantile-tail floor gives
\[
 \tau_\nu(t)=\int\min(1,t/d)\,\nu(dd)\le\mathcal K_{\mu,\rho}(N).
 \tag{AT10}
\]
For $x\ge1$, $1+\log_2x\le2x$. It follows that $w_J(d)/N\le2048\min(1,t/d)$
and $\min(1,16(t/d)^2)\le4\min(1,t/d)$. On $\nu$ their summed integral is
at most $2052\tau_\nu(t)$. On $\sigma$, using $a\le d\le1/2$, it is at most
$[1024(J+1)+16]t^2/a^2\le2048(J+1)t^2/a^2$.

Choose a finite $J_0$ after which
\[
 2048(J+1)2^{-J}/a^2\le\sqrt\rho.
 \tag{AT11}
\]
This also ensures $t\le\sqrt\rho/2$, where the elementary objective floor,
after the infimum and cap, gives $\mathcal K_{\mu,\rho}(N)\ge2\sqrt\rho\,t$.
Consequently
\[
 u_J+W_J/N\le(2052+1/2)\mathcal K_{\mu,\rho}(N)
 \le4096\mathcal K_{\mu,\rho}(N).
 \tag{AT12}
\]
Multiplying by 128 in (AT3) proves (AT4). The core's logarithmic work is
absorbed after a law-dependent onset; the tail is charged directly to (AT10).

For simultaneous certificates at all $J\ge1$, use fresh streams and
$\delta_J=\delta/[J(J+1)]$. Their sum is $\delta$, and total work through
scale $J$ is $O(4^J[\log(1/\delta)+\log(J+1)])$. On one event all certificates
are valid and (AT4) holds past $J_0$. These concern different tail thresholds,
not nested confidence intervals for the positive fraction.

\paragraph{Execution and verification boundary.}
The implementation uses integer thresholds, exact rational endpoints and
$m=8\lceil\log_2(2/\delta)\rceil+1$. A resumable probe preserves its current
coin and charges every flip; the median driver validates full authorization
before acquisition. Its tests cover pause points, reflection, threshold ties,
partial work, finite exact probabilities and failure paths. Neither Python
execution nor the new probability theorem is Lean-verified. Existing Lean
results support (AT10) and the elementary objective floor used above.
